\begin{figure*}[t]
\centering
% problem_id: 4010aaeb-28f7-5704-bc35-a5286b7efa1f (Task_3_18, ground truth ['ts'])
\begin{lstlisting}[basicstyle=\ttfamily\footnotesize, frame=single, breaklines=true, upquote=true]
(a) JSON -- database span:
{'traceId': 'bBniqPX+0EjLeW1N4kZczA==', 'spanId': '2', 'parentSpanId': '1', 'name': 'SELECT ts.orders', ..., 'db.operation': 'SELECT', 'db.sql.table': 'orders', 'db.name': 'ts', 'db.system': 'mysql'}
Llama-3.1 answer: ['ts']   (correct)

(b) CSV -- complete trace (header and both spans):
traceId,spanId,parentSpanId,name,...,db.operation,db.sql.table,db.name,db.system
bBniqPX+0EjLeW1N4kZczA==,1,-1,/api/v1/orderservice/order/refresh,...,,,,
bBniqPX+0EjLeW1N4kZczA==,2,1,SELECT ts.orders,...,SELECT,orders,ts,mysql
Llama-3.1 answer: ['mysql']   (incorrect)
\end{lstlisting}
\caption{RQ1: The same trace in JSON and CSV for Task 3.18, ``Which database is accessed most frequently in the Input Trace?'' (correct answer \texttt{ts}), with the answers of Llama-3.1. Part (a) shows the span that accesses the database, part (b) the complete trace. Middle attributes are elided (\texttt{...}).}
\label{fig:rq1_example}
\end{figure*}
